% 1-3 题 图 1-2：一系列光滑斜面（顶端与底端水平距离同为 l，倾角 φ 连续取值）
\begin{tikzpicture}[scale=1.15, font=\small]
  % 顶点 A 与底端竖直参考线
  \coordinate (A) at (0,3);
  \coordinate (B1) at (2.6,2.45);
  \coordinate (B2) at (2.6,1.65);
  \coordinate (B3) at (2.6,0.85);
  \coordinate (B4) at (2.6,0.2);
  \coordinate (D) at (2.6,3);
  % 底端线（虚线，表示底端在同一竖直线上）
  \draw[dashed] (D) -- (2.6,0.05);
  % 水平参考线
  \draw[gray!70] (A) -- (D);
  % 斜面（倾角依次增大）
  \draw[thick] (A) -- (B1);
  \draw[thick] (A) -- (B2);
  \draw[thick] (A) -- (B3);
  \draw[thick] (A) -- (B4);
  % 小球（在 B2 斜面上）
  \fill (0.95,2.22) circle (2.2pt);
  % 水平距离 l
  \draw[<->, >=Stealth] (0,-0.15) -- (2.6,-0.15) node[midway,below] {$l$};
  % 角度标注：φ1 φ2 φ3（斜面与水平线夹角）
  \pic[draw, angle radius=9mm, angle eccentricity=1.35, "$\varphi_1$"] {angle = D--A--B1};
  \pic[draw, angle radius=15mm, angle eccentricity=1.25, "$\varphi_2$"] {angle = D--A--B2};
  \pic[draw, angle radius=21mm, angle eccentricity=1.18, "$\varphi_3$"] {angle = D--A--B3};
  \pic[draw, angle radius=27mm, angle eccentricity=1.15, "$\varphi_4$"] {angle = D--A--B4};
  % 顶点标记
  \fill (A) circle (1.8pt) node[above right] {$A$};
\end{tikzpicture}
